% Problem record op_0c17d9fef967858d. This TeX file is derived from
% database/problems_json/op_0c17d9fef967858d.json by scripts/sync-tex.mjs; edit the JSON record
% and rerun the script rather than editing this file.
\section{Quantum query complexity of Triangle Finding}
\label{sec:op_0c17d9fef967858d}

\paragraph{Problem.}
What is the bounded-error quantum query complexity of finding a triangle
in an $n$-vertex graph given oracle access to its adjacency matrix?

Let $G=(V,E)$ be a simple undirected graph with $V=[n]$, and let
$A\in\{0,1\}^{n\times n}$ be its adjacency matrix.  The input is accessed
through a quantum oracle acting as
\begin{equation}
  O_A\lvert u,v,z\rangle
  =\lvert u,v,z\oplus A_{uv}\rangle .
  \label{eq:0c17-triangle-oracle}
\end{equation}
The task is to output three distinct vertices $u,v,w\in[n]$ satisfying
\begin{equation}
  A_{uv}=A_{vw}=A_{wu}=1,
  \label{eq:0c17-triangle-certificate}
\end{equation}
if such vertices exist, and otherwise report that the graph is
triangle-free.  Let $Q_{\triangle}(n)$ denote the minimum number of queries
to the oracle in \eqref{eq:0c17-triangle-oracle} required by a
bounded-error quantum algorithm for this task.

The best known general bounds are
\begin{equation}
  Q_{\triangle}(n)=O\left(n^{5/4}\right)
  \qquad\text{and}\qquad
  Q_{\triangle}(n)=\Omega(n).
  \label{eq:0c17-triangle-gap}
\end{equation}
The open question is to determine the asymptotic behavior of
$Q_{\triangle}(n)$, equivalently to improve either the upper or lower bound
in \eqref{eq:0c17-triangle-gap} until the bounds match.

\subsection*{Status}

Unsolved

\subsection*{Source}

Li and Li explicitly identify the gap between the $O(n^{5/4})$ upper
bound and the $\Omega(n)$ lower bound as open for the ordinary
adjacency-matrix problem
\sourcecite{ref:0c17-li-li}{LL25}.  The precise statement above is
contributor wording that packages this known gap as a standalone open
problem; it is not a verbatim question attributed to Li and Li or to the
authors of the earlier algorithms.

\subsection*{Progress}

\begin{itemize}
  \item Le Gall obtained a bounded-error quantum algorithm with query complexity
$\widetilde O(n^{5/4})$, using combinatorial properties specific to
unweighted triangle finding together with quantum-search and quantum-walk
techniques
\sourcecite{ref:0c17-le-gall}{LG14}.

  \item Carette, Lauri\`ere, and Magniez introduced extended learning graphs,
removing the polylogarithmic overhead from the above bound and establishing
the current best general upper bound $O(n^{5/4})$; their framework also
yields improved bounds in several sparse-graph regimes
\sourcecite{ref:0c17-carette}{CLM17}.

  \item Belovs and Rosmanis proved an $\Omega(n^{9/7}/\sqrt{\log n})$ quantum
query lower bound for the related edge-weighted triangle-sum problem and
showed that the $O(n^{9/7})$ non-adaptive learning-graph approach to
Triangle Finding is essentially optimal within that framework
\sourcecite{ref:0c17-belovs}{BR14}.
\end{itemize}

\subsection*{References}

\begin{enumerate}[label={},leftmargin=4.8em,labelsep=0.5em]
  \footnotesize
  \item[\textup{[LG14]}]\label{ref:0c17-le-gall}
  Fran\c{c}ois Le Gall, ``Improved Quantum Algorithm for Triangle Finding
via Combinatorial Arguments,'' in \emph{2014 IEEE 55th Annual Symposium on
Foundations of Computer Science (FOCS)}, 216--225 (2014).
\href{https://doi.org/10.1109/FOCS.2014.31}{doi:10.1109/FOCS.2014.31};
\href{https://arxiv.org/abs/1407.0085}{arXiv:1407.0085}.

  \item[\textup{[CLM17]}]\label{ref:0c17-carette}
  Titouan Carette, Mathieu Lauri\`ere, and Fr\'ed\'eric Magniez,
``Extended Learning Graphs for Triangle Finding,'' in \emph{34th Symposium on
Theoretical Aspects of Computer Science (STACS 2017)}, LIPIcs
\textbf{66}, Article 20 (2017).
\href{https://doi.org/10.4230/LIPIcs.STACS.2017.20}{doi:10.4230/LIPIcs.STACS.2017.20};
\href{https://arxiv.org/abs/1609.07786}{arXiv:1609.07786}.

  \item[\textup{[BR14]}]\label{ref:0c17-belovs}
  Aleksandrs Belovs and Ansis Rosmanis, ``On the Power of Non-Adaptive
Learning Graphs,'' \emph{Computational Complexity} \textbf{23}, 323--354
(2014).
\href{https://doi.org/10.1007/s00037-014-0084-1}{doi:10.1007/s00037-014-0084-1};
\href{https://arxiv.org/abs/1210.3279}{arXiv:1210.3279}.

  \item[\textup{[LL25]}]\label{ref:0c17-li-li}
  Guanzhong Li and Lvzhou Li, ``Derandomization of Quantum Algorithm for
Triangle Finding,'' \emph{Information and Computation} \textbf{304}, 105295
(2025).
\href{https://doi.org/10.1016/j.ic.2025.105295}{doi:10.1016/j.ic.2025.105295};
\href{https://arxiv.org/abs/2309.13268}{arXiv:2309.13268}.
\end{enumerate}

\subsection*{Comment}

The open question is to determine the asymptotic behavior of
$Q_{\triangle}(n)$: either an algorithm using $O(n^{5/4-\varepsilon})$
queries for some constant $\varepsilon>0$, or a lower bound of
$\Omega(n^{1+\varepsilon})$ for unrestricted bounded-error quantum
algorithms; neither is known for general unweighted graphs in the
adjacency-matrix oracle model.  The distinction between unweighted Triangle
Finding and edge-weighted Triangle Sum is essential: the stronger lower
bounds known for Triangle Sum (Progress above) must not be recorded as lower
bounds for this problem.  Li and Li, who identify the gap as open, study a
different question, namely derandomizing related triangle-finding algorithms
under an additional uniqueness promise, and do not close this gap
\sourcecite{ref:0c17-li-li}{LL25}.

\subsection*{Field}

Quantum algorithm

\subsection*{Topic}

Computational complexity and computability

\subsection*{ID}

\texttt{op\_0c17d9fef967858d}
